\section{Charakteristiky variability}\label{sec:ai-charakteristiky-variability}

Dva statistické soubory mohou mít stejný průměr i~medián, ale přesto se výrazně lišit. Hodnoty prvního mohou být soustředěny těsně kolem středu, zatímco hodnoty druhého mohou být rozptýleny po širokém intervalu. Charakteristiky variability popisují právě tuto rozptýlenost.

\begin{definition}\label{def:ai-rozptyl}
    \emph{Rozptyl} statistického souboru
    \(\mathbf{X}=(x_1,\ldots,x_n)\) definujeme vztahem
    \[
        \Var(\mathbf{X})
        =
        \sigma_{\mathbf{X}}^2
        =
        \frac1n\sum_{i=1}^{n}
        (x_i-\average{\mathbf{X}})^2.
    \]
    \emph{Směrodatná odchylka} tohoto souboru je číslo
    \[
        \sigma_{\mathbf{X}}
        =
        \sqrt{\Var(\mathbf{X})}.
    \]
\end{definition}

Rozptyl je vždy nezáporný a~je nulový právě tehdy, když jsou všechny hodnoty stejné. Protože pracuje se čtverci odchylek, jeho jednotka je druhou mocninou původní jednotky. Směrodatná odchylka tento problém odstraňuje a~vyjadřuje typickou velikost odchylky ve stejné jednotce jako původní data.

\begin{figure}[ht]
    \centering
    \begin{tikzpicture}[x=0.9cm,y=1cm,every node/.style={font=\normalsize}]
    \draw[diredge] (0,1.4) -- (10.5,1.4);
    \draw[diredge] (0,-0.3) -- (10.5,-0.3);
    \node[left] at (0,1.4) {\(\mathbf{X}\)};
    \node[left] at (0,-0.3) {\(\mathbf{Y}\)};
    \foreach \x in {4.3,4.7,5,5.2,5.8}{
        \fill[blue!75!black] (\x,1.4) circle (2.5pt);
    }
    \foreach \x in {1,3,5,7,9}{
        \fill[orange!85!black] (\x,-0.3) circle (2.5pt);
    }
    \draw[red!75!black,very thick] (5,1.05) -- (5,1.75);
    \draw[red!75!black,very thick] (5,-0.65) -- (5,0.05);
    \node[right,red!75!black] at (5.1,1.75)
        {\(\average{\mathbf{X}}=\average{\mathbf{Y}}\)};
    \node[below] at (5,-0.75)
        {\(\Var(\mathbf{X})<\Var(\mathbf{Y})\)};
\end{tikzpicture}

    \caption{Soubory se stejným průměrem a~různou variabilitou.}
    \label{fig:ai-variabilita}
\end{figure}

\subsection{Cauchyova--Schwarzova nerovnost}

Při práci s~rozptylem a~korelací budeme potřebovat následující základní nerovnost.

\begin{theorem}[Cauchyova--Schwarzova nerovnost]\label{thm:ai-cauchy-schwarz}
    Pro libovolná reálná čísla
    \(a_1,\ldots,a_n,b_1,\ldots,b_n\) platí
    \[
        \left|
            \sum_{i=1}^{n}a_ib_i
        \right|
        \leqslant
        \sqrt{\sum_{i=1}^{n}a_i^2}
        \sqrt{\sum_{i=1}^{n}b_i^2}.
    \]
    Rovnost nastává právě tehdy, když existuje \(\lambda\in\R\) takové, že
    \(b_i=\lambda a_i\) pro všechna \(i\), nebo když jsou všechna \(a_i\) nulová.
\end{theorem}

Důkaz věty zde nebudeme uvádět. Geometricky jde o~zobecnění skutečnosti, že absolutní hodnota skalárního součinu dvou vektorů nepřekročí součin jejich délek.

\subsection{Vztah průměru, mediánu a~směrodatné odchylky}

\begin{proposition}\label{prop:ai-prumer-median-sigma}
    Pro každý statistický soubor
    \(\mathbf{X}=(x_1,\ldots,x_n)\) platí
    \[
        \abs{\average{\mathbf{X}}-\median(\mathbf{X})}
        \leqslant
        \sigma_{\mathbf{X}}.
    \]
\end{proposition}

\begin{proof}
    Označme \(\mu=\average{\mathbf{X}}\) a~\(m=\median(\mathbf{X})\).
    Stačí vyřešit případ \(\mu\geqslant m\); případ \(\mu<m\) dostaneme stejným postupem po změně znamének.

    Položme
    \[
        A=\set{i:x_i\leqslant m},
        \qquad
        k=\sizeof{A}.
    \]
    Z~definice mediánu plyne \(k\geqslant n/2\). Pro \(i\in A\) definujme
    \(a_i=\mu-x_i\). Potom \(a_i\geqslant\mu-m\). Označíme-li
    \[
        S=\sum_{i\in A}a_i,
    \]
    dostáváme \(S\geqslant k(\mu-m)\). Součet všech odchylek od průměru je nulový, a~proto
    \[
        \sum_{i\notin A}(x_i-\mu)=S.
    \]

    Z~Cauchyovy--Schwarzovy nerovnosti použité zvlášť na obě skupiny plyne
    \[
        \sum_{i\in A}(x_i-\mu)^2
        \geqslant
        \frac{S^2}{k}
        \quad\text{a}\quad
        \sum_{i\notin A}(x_i-\mu)^2
        \geqslant
        \frac{S^2}{n-k}.
    \]
    Je-li \(k=n\), jsou všechny hodnoty nejvýše \(\mu\), a~nulový součet odchylek vynutí \(x_i=\mu=m\); tvrzení je okamžité. Pro \(k<n\) sečtením nerovností získáme
    \begin{align*}
        \sigma_{\mathbf{X}}^2
        &=
        \frac1n\sum_{i=1}^{n}(x_i-\mu)^2\\
        &\geqslant
        \frac1n
        \left(
            \frac{S^2}{k}+\frac{S^2}{n-k}
        \right)
        =
        \frac{S^2}{k(n-k)}\\
        &\geqslant
        \frac{k}{n-k}(\mu-m)^2
        \geqslant
        (\mu-m)^2.
    \end{align*}
    Poslední nerovnost plyne z~\(k\geqslant n/2\). Po odmocnění tedy
    \(\mu-m\leqslant\sigma_{\mathbf{X}}\).
\end{proof}

\importantbox{Tvrzení neříká, že průměr a~medián musí být blízko v~absolutním měřítku. Říká, že jejich vzdálenost nemůže překročit směrodatnou odchylku daného souboru.}
